\chapter{Nachweise, Grenzen und Rückverfolgbarkeit}\label{chap:evidence}
\section{Was ein Nachweis aussagt}
Code-Verifikation prüft implementierte Gleichungen, Lösungsverifikation ihre
räumliche/zeitliche Auflösung. Validierung bewertet Modellabweichung gegenüber
geeigneten unabhängigen Messdaten. Dieser Theorieband beansprucht die ersten
beiden Nachweisarten nur für die angegebenen Referenzen, keine experimentelle
Validierung. Manufactured Solutions setzen bekannte Felder in die starke Form
ein und leiten Quellen/Randwerte unabhängig ab; der Sandia-Bericht
\cite{sandia-mms} erläutert dieses Prinzip. Gleicher Output zweier Solver allein
ist keine vollständige Verifikation.

Jedes Kapitel enthält Annahmen, Variablen, Einheiten, Ableitung, implementierte
Form und Nutzungsgrenzen. \path{Doc/latex/THEORY_VERIFICATION.md} protokolliert
die tatsächliche Nachrechnung und Wiederholung der kleinen Referenztests.
\path{Doc/latex/theory-map.json} ordnet je Bereich Quellkapitel,
Implementierung, kanonische Verträge und analytische Referenzen, Patch-, Bilanz- und
Konvergenztests zu. Die übergreifende Quellenmatrix bleibt
\path{Doc/latex/coverage.json}. Der dortige Original-Status beschreibt die
Herkunft des Initialstands; \texttt{theory\_\allowbreak coverage} beschreibt separat den
Ausarbeitungsstand dieses Bandes und ist keine automatische Freigabe.

\section{Nachweisübersicht nach Verfahren}
\begin{longtable}{p{0.22\textwidth}p{0.69\textwidth}}
\caption{Unabhängige Referenzen und zugehörige bestehende Tests.}\label{tab:evidence}\\
\toprule Bereich & Referenz und Prüfquelle\\\midrule\endfirsthead
\toprule Bereich & Referenz und Prüfquelle\\\midrule\endhead
Konventionen & Achsrotation, Energieinvarianz, Strain-/Stress-Charge-Konversion;
\path{tests/unit/test_piezo_material.py}.\\
Tet4/Tet10 & Simplexmomente, alle Eckpermutationen, Monomreproduktion;
\path{tests/unit/test_tet10_geometry.py} und
\path{tests/unit/test_tet10_operators.py}.\\
Wärme & Affiner Fluss, Robin-Quader, Quellen-/Reservoirbilanz;
\path{tests/patch/test_heat_patch.py}; vier Netze in
\path{tests/convergence/test_heat_convergence.py}.\\
Elektrostatik & Kondensator $C=\kappa A/L$, D-Flussvorzeichen, Energie;
\path{tests/patch/test_electrostatics_patch.py}; vier Netze in
\path{tests/convergence/test_electrostatics_convergence.py}.\\
Elastizität & Starrkörper-/Affinpatch, unabhängiger Zugzustand, Kraft/Moment;
\path{tests/patch/test_elasticity_patch.py} und
\path{tests/convergence/test_elasticity_convergence.py}.\\
Thermomechanik & Freie/behinderte gedrehte Expansion, gekoppelte MMS;
\path{tests/patch/test_thermomechanics_patch.py} und
\path{tests/convergence/test_thermomechanics_convergence.py}.\\
Piezo & Affine Enthalpievariation, d33-Reziprozität und zweite unabhängige
1D-FE-Rechnung; \path{tests/regression/test_piezo_d33_benchmark.py},
\path{tests/reference/piezo_bar_1d.py}.\\
Thermopiezo & Geklemmter thermischer Potentialgrenzfall, lineare Überlagerung;
\path{tests/unit/test_thermopiezoelectric.py}; Material-/Elektrodenmodell auf
vier Netzen in \path{tests/convergence/test_thermopiezo_model_convergence.py}.\\
Elektroden & Drei analytische Gesamtladungs-/Potentialfälle;
\path{tests/patch/test_electrodes.py}.\\
Masse/Wärme & Exakte Momente, positive Masse und HRZ-Gesamterhaltung;
\path{tests/unit/test_storage.py} und
\path{tests/unit/test_tet10_operators.py}.\\
Backward Euler & Analytische Abkühlrate, Quellensprung, Speicher-/Randbilanz;
\path{tests/unit/test_transient_heat.py} und
\path{tests/convergence/test_transient_heat_convergence.py}.\\
Modal & Unabhängiger Ein-DOF-Oszillator und offener/kurzer d33-Stab;
\path{tests/unit/test_piezo_modal.py} und
\path{tests/convergence/test_piezo_modal_convergence.py}.\\
Newmark & Energie und zweite Zeitordnung auf festem Raummodell, bewegte
Ränder/Trägheit; \path{tests/unit/test_piezo_newmark.py}.\\
Harmonik & Komplexer Oszillator, DC, Resonanzfehler und Newmark-Parität;
\path{tests/unit/test_piezo_harmonic.py}.\\
Tet10-Solver & Vier-Netze-Feldnormen je Physik, integrierte Energie und
Bilanzen; \path{Doc/55_TET10_ACCEPTANCE.md} bündelt die Einzelnachweise.\\
Tet10-Dynamik & Eigene Stab-Modal-/Harmonik-Raumverfeinerung und Newmark-
Zeitprüfung; \path{tests/convergence/test_tet10_modal.py},
\path{tests/convergence/test_tet10_harmonic_convergence.py},
\path{tests/unit/test_tet10_newmark.py}.\\
PETSc & Native real/complex-Parität, FieldSplit-Abbruch, freie Feldresiduen;
\path{tests/petsc}. Fehlende Laufzeit ist lokal sichtbar, in nativer CI ein Fehler.\\
Indikator & Analytische Zellen-/Flussterme und getrennte Fehler-/Indikator-
Netzstudie; \path{tests/unit/test_heat_indicator.py},
\path{tests/convergence/test_heat_indicator_convergence.py}.\\
Surrogate & Unabhängige lineare/quadratische Funktionen, Ridge-Intercept,
Holdouts, Box-/Rangfehler; \path{tests/unit/test_surrogates.py} und
\path{tests/regression/test_surrogate_example.py}.\\
\bottomrule
\end{longtable}

\section{Wie die Referenztests interpretiert werden}
Algebraische Identitäten und exakt darstellbare Zuordnungen können exakt
verglichen werden. Nach linearen Lösevorgängen werden relative und absolute
Toleranzen aus SI-Skalen und Kondition vorab festgelegt. Beispielsweise gelten
$10^{-12}$ V/C im synthetischen Elektrodenquader, aber $2\cdot10^{-22}$ C
für einen Kondensator mit Ladung in der Größenordnung $10^{-9}$ C.
Ein gemeinsamer dimensionsloser absoluter Grenzwert wäre fachlich falsch.
Die hier wiederholten Tests behalten ihre bestehenden Toleranzen.

Die historischen vierstufigen Nachweise aus Doc/55 und Doc/25 bleiben anwendbar,
weil dieser PR keine numerischen Operatoren oder Materialien ändert.
P1-Fenster für glatte Probleme sind typischerweise L2 1,8--2,2 und H1
0,85--1,15, P2 2,7--3,3 und 1,8--2,2 auf den letzten Verfeinerungen.
Diese Fenster werden nicht nachträglich angepasst. Exportmasken, Schwerpunkt-
orte und Energieintegrale sind eigene Verträge und keine Konvergenzordnungen.

\section{Ein minimaler nachvollziehbarer Matrixeinstieg}
Die folgende lokale Rechnung reproduziert den Patch aus
Gleichung~\ref{eq:piezo-patch-example}, ohne eine globale Lagerung oder einen
Solver als Referenz zu benötigen. $x^T A x/2$ ist Enthalpie, keine positive Norm.
Alle Zahlen sind synthetische SI-Daten.
\begin{lstlisting}[language=Python,caption={Lokale Enthalpie des affinen Piezo-Patches.},label={code:piezo-patch}]
import numpy as np
from fem_toolbox.elements import Tet4Geometry
from fem_toolbox.elements.piezoelectric import piezoelectric_matrix
from fem_toolbox.materials import PiezoelectricMaterial

points = np.array([[0., 0., 0.], [1., 0., 0.],
                   [0., 1., 0.], [0., 0., 1.]])
e = np.zeros((3, 6)); e[0, 0] = 2.
material = PiezoelectricMaterial(
    np.diag([10., 20., 30., 4., 5., 6.]), e,
    np.diag([3., 4., 5.]))
A = piezoelectric_matrix(Tet4Geometry(points), material)
u = np.zeros((4, 3)); u[:, 0] = 0.2 * points[:, 0]
x = np.r_[u.ravel(), 0.3 * points[:, 0]]
enthalpy = float(x @ A @ x / 2)  # 37/1200 J
\end{lstlisting}

\section{Checkliste zur Anwendung}
Vor einer Rechnung müssen Geometrie/Orientierung, Materialachsen/Einheiten,
Eindeutigkeit, Randarten und Ereignis-/Frequenzraster zum Modell passen.
Nach einer Rechnung werden freie Feldresiduen, physikalische Bilanzen,
rohe Feldherkunft und getrennte Raum-/Zeitfehler bewertet.
Kostenvergleiche sind nur im gemessenen Raster gültig.
Eine später erhobene Messabweichung darf weder als Diskretisierungsfehler noch
als empirischer Surrogatfehler umetikettiert werden.

Der aktuelle Band liefert keine fachliche Fremdfreigabe. Eigenprüfung und
CI sind dokumentiert; fehlendes unabhängiges Review wird im PR benannt.
Der Anwendungsband ist als Kandidat in \#208 ausgearbeitet.
API- und V\&V-Bände bleiben Gegenstand von \#209--\#210,
Gesamtabnahme \#211 und Release \#66. Das abschließende Handbuch kann erst nach
diesen Gates als vollständig bezeichnet werden.
